% part: intuitionistic-logic
% chapter: propositions-as-types
% section: type-preservation

\documentclass[../../../include/open-logic-section]{subfiles}

\begin{document}

\olfileid[id]{int}{pty}{tpr}

\olsection{Preservasi Tipe}

Konversi reduksi dan permutasi yang telah dibahas juga dapat didefinisikan
langsung bagi sistem \Log{tN2i} dan \Log{tN3i}. Sistem \Log{tN3i}
mendefinisikan tipe suatu term-$\lambd$~$N$ relatif terhadap konteks~$\Gamma$
sebagai !!{formula}~$!A$ sedemikian sehingga $\Gamma \Sequent N:!A$
mempunyai !!a{proof} dalam \Log{tN3i}. Fakta bahwa aturan reduksi-$\beta$
bagi term-$\lambd$ mencerminkan secara tepat konversi reduksi dan permutasi
pada !!{proof}s mempunyai akibat penting: tipe suatu term-$\lambd$ relatif
terhadap konteks~$\Gamma$ tetap sama sesudah konversi-$\beta$. Sifat ini
disebut \emph{preservasi tipe}.

\begin{prop}
Jika $N$ bertipe~$!A$ relatif terhadap $\Gamma$ dan $N \redone N'$, maka
$N'$ bertipe~$!A$ relatif terhadap~$\Gamma$.
\end{prop}

\begin{proof}
  Melalui induksi pada~$N$ dan tinggi !!{proof}s~$\delta$ bagi
  $\Gamma \Sequent N:!A$, kita dapat dengan mudah menetapkan hal berikut.
  Jika $M$ merupakan subterm dari~$N$, maka $\delta$ memuat
  sub-!!{proof}~$\delta_1$ yang berakhir dengan
  $\Gamma' \Sequent M:!B$. Jika $M$ merupakan redex, suatu konversi dapat
  diterapkan pada~$\delta_1$. Dengan menerapkan reduksi pada~$\delta_1$, kita memperoleh
  !!a{proof}~$\delta_1'$ bagi $\Gamma' \Sequent M':!B$. Dari pemeriksaan
  konversi-konversi \Log{tN3i}, $M'$ adalah kontraktum dari~$M$. Mengganti
  $\delta_1$ dalam~$\delta$ dengan $\delta_1'$ menghasilkan
  !!a{proof}~$\delta'$ bagi $\Gamma \Sequent N':!A$, dengan $N'$ diperoleh
  dari~$N$ dengan mengganti subterm~$M$ oleh~$M'$.

  Sebagai contoh, tinjau konversi reduksi bagi \Intro{\land} yang diikuti
  \Elim{\land}:
  \[
  \bottomAlignProof
  \AxiomC{}
  \Deduce$\Gamma \fCenter N_1 : !B_1$
  \AxiomC{}
  \Deduce$\Gamma \fCenter N_2 : !B_2$
  \RightLabel{\Intro{\land}}
  \BinaryInf$\Gamma \fCenter \pair{N_1}{N_2} : !B_1 \land !B_2$
  \RightLabel{\Elim{\land}[i]}
  \UnaryInf$\Gamma \fCenter \proj{i}{\pair{N_1}{N_2}} : !B_i$
  \DisplayProof
  \quad\redone\quad
  \bottomAlignProof
  \AxiomC{}
  \Deduce$\Gamma \fCenter N_i : !B_i$
  \DisplayProof
  \]
  Kita melihat bahwa term bukti dari !!{proof} di sebelah kanan, yakni
  $\delta_1'$, adalah~$N_i$. Term ini tepat merupakan hasil penerapan
  konversi-$\beta$ pada term bukti $\proj{i}{\pair{N_1}{N_2}}$ dari
  !!{proof} di sebelah kiri~($\delta_1$).

  Atau, tinjau inferensi \Intro{\lif} yang diikuti \Elim{\lif}, beserta
  konversi reduksi yang diterapkan kepadanya:
  \begin{multline*}
  \bottomAlignProof
  \AxiomC{}
  \RightLabel{$\delta_2$}
  \Deduce$x: !B_1, \Gamma \fCenter N : !B_2$
  \RightLabel{\Intro{\lif}}
  \UnaryInf$\Gamma \fCenter \lambd[\typeof{x}{!B_1}][N] : !B_1 \lif !B_2$
  \AxiomC{}
  \RightLabel{$\delta_3$}
  \Deduce$\Gamma \fCenter M : !B_1$
  \RightLabel{\Elim{\lif}}
  \BinaryInf$\Gamma \fCenter (\lambd[\typeof{x}{!B_1}][N] M) : !B_2$
  \DisplayProof
  \quad\redone\\
  \bottomAlignProof
  \AxiomC{}
  \RightLabel{$\delta_3$}
  \Deduce$\Gamma \fCenter M : !B_1$
  \RightLabel{$\Subst{\delta_2}{\delta_3}{x:!B_1}$}
  \Deduce$\Gamma \fCenter \Subst{N}{M}{x} : !B_2$
  \DisplayProof
  \end{multline*}
  Sekali lagi, term bukti dari !!{proof} di sebelah kanan tepat merupakan
  hasil reduksi-$\beta$ atas term bukti dari !!{proof} di sebelah kiri.
  Tentu saja, kita harus memverifikasi dengan cermat (melalui induksi pada
  tinggi~$\delta_2$) bahwa mengganti semua aksioma berbentuk
  $\Gamma' \Sequent x:!B_1$ dalam~$\delta_2$ dengan
  !!{proof}~$\delta_3$ bagi $\Gamma \Sequent M:!B_1$ benar-benar
  menghasilkan !!a{proof} $\Subst{\delta_2}{\delta_3}{x:!B_1}$ bagi
  $\Gamma \Sequent \Subst{N}{M}{x}:!B_2$.
\end{proof}

Korespondensi erat antara !!{proof}s dan term-$\lambd$ mempunyai akibat
lain. Jika $N$ merupakan term-$\lambd$ bertipe $!A$ relatif terhadap
konteks~$\Gamma$ dan memuat redex~$M$ (yakni, $N$ tidak berada dalam bentuk
normal), maka bukti \Log{tN3i}~$\delta$ yang bersesuaian bagi
$\Gamma \Sequent N:!A$ memuat sub-!!{proof} (yang berakhir dengan
$\Gamma' \Sequent M:!B$) tempat suatu konversi dapat diterapkan. Jika
$N \redone N'$ dengan mengganti $M$ oleh kontraktumnya, hasil yang
bersesuaian dari penerapan reduksi pada~$\delta$ adalah !!a{proof} bagi
$\Gamma \Sequent N':!A$. Jadi, setiap term-$\lambd$ yang tidak berada dalam
bentuk normal tereduksi-$\beta$ menjadi term lain. Sifat ini disebut
\emph{kemajuan}.

Kemajuan bersama preservasi tipe memastikan bahwa reduksi-$\beta$ atas
term-$\lambd$ tidak pernah ``macet'': jika suatu term bertipe baik dan tidak
berada dalam bentuk normal, term itu tereduksi menjadi term bertipe baik lain
(dan term tersebut bertipe sama).

\end{document}
